
\documentclass[10pt,aspectratio=169]{beamer}

% ------------------------------------------------------------
% Packages
% ------------------------------------------------------------
\usepackage[T1]{fontenc}
\usepackage[utf8]{inputenc}
\usepackage{lmodern}
\usepackage{microtype}
\usepackage{amsmath,amssymb}
\usepackage{graphicx}
\usepackage{booktabs}
\usepackage{xcolor}
\usepackage{listings}

% ------------------------------------------------------------
% Theme
% ------------------------------------------------------------
\usetheme{Madrid}
\usecolortheme{default}
\setbeamertemplate{navigation symbols}{}

\definecolor{ustcblue}{RGB}{0,55,120}
\definecolor{softgray}{RGB}{245,246,248}
\definecolor{darkred}{RGB}{140,25,25}

\setbeamercolor{structure}{fg=ustcblue}
\setbeamercolor{frametitle}{fg=white,bg=ustcblue}
\setbeamercolor{title}{fg=white,bg=ustcblue}
\setbeamercolor{block title}{fg=white,bg=ustcblue}
\setbeamercolor{block body}{bg=softgray,fg=black}

\lstset{
    basicstyle=\ttfamily\scriptsize,
    backgroundcolor=\color{softgray},
    frame=single,
    breaklines=true,
    columns=fullflexible,
    keepspaces=true,
    showstringspaces=false
}

% ------------------------------------------------------------
% Metadata
% ------------------------------------------------------------
\title[FRC Weekly Update]{FRC Project Weekly Update}
\subtitle{Formation Baseline Closed; GS/MRR Takeover Started}
\author{Juan Pablo Solís Ruiz}
\institute{USTC --- Supervisor: Prof. Ren Haijun}
\date{October 13, 2026}

% ------------------------------------------------------------
\begin{document}

% ------------------------------------------------------------
\begin{frame}
    \titlepage
\end{frame}

% ------------------------------------------------------------
\begin{frame}{Status Since Last Update}
    \begin{block}{Main message}
        The validated formation workflow remains closed. The main new progress
        was a preliminary forensic audit of the inherited GS/MRR initializer,
        not a new production simulation.
    \end{block}

    \vspace{0.3cm}

    \begin{itemize}
        \item Formation baseline reproduction remains validated.
        \item The \(4~\mu\mathrm{s}\) formation extension was documented as a
        controlled exploratory run.
        \item The inherited \texttt{hybrid\_init.py} / MRR-GS workflow was
        reviewed conceptually.
        \item No GS/MRR physics or tilt perturbation has been modified yet.
        \item Progress was limited by a temporary high academic and project load.
    \end{itemize}
\end{frame}

% ------------------------------------------------------------
\begin{frame}{Formation Work Is Now a Reference, Not the Main Focus}
    \begin{columns}[T,onlytextwidth]
        \begin{column}{0.48\textwidth}
            \centering
            \includegraphics[width=\linewidth]{img/framework_v_01/Bz_center.png}

            \vspace{0.1cm}
            \small Validated 2000-step baseline regression
        \end{column}

        \begin{column}{0.48\textwidth}
            \centering
            \includegraphics[width=\linewidth]{img/formation_4us/Bz_radial_2_3_4us.png}

            \vspace{0.1cm}
            \small Exploratory \(2\rightarrow4~\mu\mathrm{s}\) extension
        \end{column}
    \end{columns}

    \vspace{0.3cm}

    \begin{block}{Conclusion}
        Formation is now frozen as a validated framework example. I do not plan
        to spend more compute on formation until new diagnostics or a specific
        physics question require it.
    \end{block}
\end{frame}

% ------------------------------------------------------------
\begin{frame}{What the \(4~\mu\mathrm{s}\) Extension Showed}
    \begin{columns}[T,onlytextwidth]
        \begin{column}{0.42\textwidth}
            \centering
            \includegraphics[width=\linewidth]{img/formation_4us/Bz_xz_step004000.png}

            \vspace{0.1cm}
            \small \(B_z(x,z)\) at \(t=4~\mu\mathrm{s}\)
        \end{column}

        \begin{column}{0.53\textwidth}
            \begin{itemize}
                \item The imposed drive reaches
                \[
                A_{\mathrm{time}}=1
                \quad\text{at}\quad
                t=4~\mu\mathrm{s}.
                \]
                \item The strongest sampled central compression occurred earlier,
                near
                \[
                t\simeq1.7~\mu\mathrm{s}.
                \]
                \item The behavior was non-monotonic:
                \[
                \text{compression}
                \rightarrow
                \text{relaxation}
                \rightarrow
                \text{renewed compression}.
                \]
                \item Field reversal persisted at the end of the ramp.
            \end{itemize}
        \end{column}
    \end{columns}

    \vspace{0.2cm}
    \small
    The \(B_z=0\) contour is still treated only as a field-reversal proxy, not a
    formal magnetic separatrix.
\end{frame}

% ------------------------------------------------------------
\begin{frame}{GS/MRR Takeover: First Forensic Pass}
    \begin{block}{Recovered workflow}
        \[
        (\beta_s,\delta_s,n,B_{\rm ext})
        \rightarrow
        \text{MRR-1 radial closure}
        \rightarrow
        \text{fixed-boundary GS}
        \rightarrow
        \texttt{init.bp}
        \rightarrow
        \text{hybrid-PIC tilt run}.
        \]
    \end{block}

    \vspace{0.25cm}

    \begin{itemize}
        \item \texttt{hybrid\_init.py} generates both the initial fields and
        the external-file particles.
        \item The tilt seed is already embedded in the particle momenta in
        \texttt{init.bp}.
        \item Therefore changing the tilt constants inside the WarpX input alone
        does not change the actual initial perturbation.
        \item The initializer appears deterministic because the particle RNG seed
        is fixed.
    \end{itemize}
\end{frame}

% ------------------------------------------------------------
\begin{frame}{MRR-1 Model Reconstructed}
    \begin{block}{Pressure closure}
        \[
        \bar p(x)
        =
        \beta_s e^{-\alpha x}
        (1+\sigma x)^n,
        \qquad
        x=\frac{\psi}{B_e R_s^2}.
        \]
    \end{block}

    \begin{columns}[T,onlytextwidth]
        \begin{column}{0.48\textwidth}
            \begin{itemize}
                \item Current defaults:
                \[
                \beta_s=0.6,\quad n=2.
                \]
                \item The code uses:
                \[
                \delta_s=0.4\delta_s^{RR}.
                \]
                \item Reconstructed values:
                \[
                \alpha\simeq31.5351,
                \quad
                \sigma\simeq8.31255.
                \]
            \end{itemize}
        \end{column}

        \begin{column}{0.48\textwidth}
            \begin{block}{GS equation}
                \[
                \Delta^\star\psi
                =
                -\mu_0 r^2 P'(\psi).
                \]
                \[
                B_z=\frac{1}{r}\frac{\partial\psi}{\partial r},
                \qquad
                B_r=-\frac{1}{r}\frac{\partial\psi}{\partial z}.
                \]
            \end{block}
        \end{column}
    \end{columns}
\end{frame}

% ------------------------------------------------------------
\begin{frame}{Important Audit Findings}
    \begin{itemize}
        \item The current initializer uses
        \[
        B_{\rm external}=0.26~\mathrm{T},
        \]
        while the reference PDF example appears to document a
        \(0.13~\mathrm{T}\) case.
        \item The MRR scale parameter
        \[
        R_s^{\rm MRR}=0.25~\mathrm{m}
        \]
        should not be confused with the geometry extracted from the converged
        equilibrium.
        \item The initializer writes
        \[
        R_0\simeq0.1707~\mathrm{m},
        \qquad
        Z_0\simeq0.3244~\mathrm{m}
        \]
        for the tilt seed.
        \item The inherited materials also contain evolved quantities such as
        \(R_s\), \(L_s\), \(E\), and \(S=R_s/d_i\), which are different objects.
    \end{itemize}
\end{frame}

% ------------------------------------------------------------
\begin{frame}{Tilt Seed: What Is Already Known}
    \begin{block}{Implemented perturbation}
        The seed is an \(n=1\) axial-velocity perturbation:
        \[
        \delta u_z
        =
        u_0
        N
        \frac{r}{R_0}
        \left(
        \frac{-\psi}{\psi_O}
        \right)^2
        \cos\theta,
        \qquad
        \psi<0.
        \]
    \end{block}

    \vspace{0.3cm}

    \begin{itemize}
        \item Seed amplitude:
        \[
        \epsilon=0.01.
        \]
        \item Effective Alfv\'en speed written by the initializer:
        \[
        V_A\simeq1.35\times10^5~\mathrm{m/s}.
        \]
        \item Peak seed velocity:
        \[
        \epsilon V_A\simeq1.35\times10^3~\mathrm{m/s}.
        \]
        \item No \(n>1\) perturbation should be added before reproducing the
        inherited \(n=1\) case.
    \end{itemize}
\end{frame}

% ------------------------------------------------------------
\begin{frame}{Workload Constraint and Current Plan}
    \begin{block}{Context}
        During the last two weeks, progress was limited by coursework,
        plasma-physics regularization, and a separate project deadline. I did
        not want to mix uncontrolled code changes with incomplete attention.
    \end{block}

    \vspace{0.3cm}

    \begin{itemize}
        \item No new GS/MRR run was launched.
        \item No tilt physics was modified.
        \item The work done was deliberately limited to reading, documenting,
        and separating definitions.
        \item This keeps the next step clean: reproduce first, then change
        physics.
    \end{itemize}
\end{frame}

% ------------------------------------------------------------
\begin{frame}{Next Concrete Steps}
    \begin{enumerate}
        \item Run \texttt{hybrid\_init.py} alone, without launching WarpX.
        \item Print and save a normalization table:
        \[
        R_0,\ Z_0,\ B_{\rm ext},\ n_0,\ V_A,\ d_i,\ S,\ E,\ S/E.
        \]
        \item Verify that generated \texttt{init.bp} matches the inherited input.
        \item Create a short WarpX smoke run for the GS/MRR case.
        \item Build the first tilt diagnostic:
        \[
        A_1(t),\qquad
        \ln A_1(t),\qquad
        \gamma_{\rm tilt}\tau_A.
        \]
    \end{enumerate}

    \vspace{0.25cm}

    \begin{center}
        \Large
        Priority: GS/MRR reproduction before parameter studies.
    \end{center}
\end{frame}

% ------------------------------------------------------------
\begin{frame}{Request for Feedback}
    \begin{block}{Decision point}
        I propose not to modify the equilibrium or add additional perturbation
        modes until the inherited \(n=1\) tilt case is reproduced.
    \end{block}

    \vspace{0.3cm}

    Feedback requested:
    \begin{itemize}
        \item Is the current priority correct: reproduce GS/MRR tilt first?
        \item Should the first reproduced case be \(B_{\rm ext}=0.26~\mathrm{T}\)
        seeded?
        \item Which diagnostic should be preferred for the tilt amplitude:
        density displacement, magnetic-axis displacement, or Fourier \(n=1\)
        amplitude?
    \end{itemize}
\end{frame}

% ------------------------------------------------------------
\end{document}
